\section{Methods}

\subsection{Dataset: African Natural Product-Derived Antimalarials}

We extracted \num{17} antimalarial candidate molecules from Project 1 Set A\cite{sao2026p1}, a curated panel enriched for African natural product (ANP) structural motifs prioritized by multi-parameter optimization (MPO $\geq 0.70$) and Resistance-Resilience Score (RRS) across four malaria drug targets (PfDHFR, PfCRT, PfATP4, PfClpP). The seed library was derived from \num{396} ANP structures sourced from ANPDB\cite{betow2025} and AfroDb\cite{ntie2017afrodb}, encompassing traditional antimalarial botanical species including \textit{Cryptolepis sanguinolenta}, \textit{Enantia chlorantha}, and \textit{Nauclea latifolia}. All molecules were represented as canonical SMILES strings and validated using RDKit\cite{rdkit}.

\subsubsection{African-Chemotype Characterization}

To quantify ANP structural identity, we computed three descriptors for each molecule:

\textbf{(i) Fraction sp\textsuperscript{3} carbons ($F_{sp^3}$):} Calculated as the ratio of sp\textsuperscript{3}-hybridized carbons to total carbons. ANPs typically exhibit $F_{sp^3} > 0.5$ due to saturated ring systems and polycyclic frameworks, contrasting with synthetic drugs ($F_{sp^3} \sim 0.3$--0.4$)\cite{lovering2009escape}.

\textbf{(ii) Stereocenter count:} Enumerated using RDKit's \texttt{FindMolChiralCenters} function after stereochemistry assignment. ANP alkaloids, triterpenoids, and macrolides frequently contain 4--10 stereocenters\cite{newman2020natural}.

\textbf{(iii) African-Chemotype Structural Index (ACSI):} Defined as the mean Tanimoto similarity (Morgan fingerprint, radius=2) to the five nearest neighbors within the 396-member ANP reference set. ACSI ranges from 0 (no ANP similarity) to 1 (perfect match); values $>0.5$ indicate strong ANP character.

Activity labels (active/inactive) were derived from docking-score thresholds applied in the upstream P1 multi-target triage workflow (see \LED{001} for label assignment criteria and validation). The final panel comprised \num{15} inactive and \num{2} active molecules (class imbalance \num{15}:\num{2}), creating a challenging benchmark representative of hit-discovery screening where active hits are rare.

\subsection{Classical Baseline: ECFP4 Fingerprints}

Extended-Connectivity Fingerprints (ECFP4)\cite{rogers2010} were generated using RDKit (\texttt{GetMorganFingerprintAsBitVect}, radius=2, 2048 bits). A support vector machine (SVM) with radial basis function (RBF) kernel was trained using scikit-learn (v1.3.0)\cite{sklearn}. Hyperparameters $C$ and $\gamma$ were set to default values (1.0 and \texttt{scale}, respectively). Performance was evaluated via leave-one-out cross-validation (LOO-CV), where each of the \num{17} molecules served as a held-out test instance while training on the remaining \num{16}. Predictions were scored using the SVM decision function, and area under the receiver operating characteristic curve (AUC-ROC) was computed as the primary metric.

This baseline represents the standard molecular fingerprinting approach used in P1 scaffold-recovery validation (\qty{69.3}{\percent} scaffold conservation)\cite{sao2026p1} and P3 quantum-inspired benchmarks (AUC \num{0.9475})\cite{sao2026p3}.

\subsection{Quantum Molecular Encoding}

\subsubsection{BondOrderMatrix Representation}

Following Boy et al.\cite{boy2025}, molecular structures were encoded as BondOrderMatrix representations preserving stereochemistry and bond topology:

\begin{equation}
\mathbf{M}_{ij} = 
\begin{cases}
0.5 \times Z_i^{2.4} & i = j \text{ (diagonal)} \\
\frac{Z_i \times Z_j}{b_{ij}} & i \neq j \text{ (off-diagonal)}
\end{cases}
\end{equation}

where $Z_i$ is the atomic number of atom $i$, and $b_{ij}$ is the bond order between atoms $i$ and $j$ (1 for single, 1.5 for aromatic, 2 for double, 3 for triple bonds).

\textbf{Stereochemical encoding:}
\begin{itemize}
    \item \textbf{Z/E isomers:} Z-configuration bonds assigned negative bond orders ($b_{ij} < 0$).
    \item \textbf{R/S chirality:} S-configuration stereocenters assigned negative diagonal elements ($\mathbf{M}_{ii} < 0$).
\end{itemize}

This encoding directly embeds 3D stereochemical information into the quantum state, unlike classical fingerprints which treat stereoisomers as distinct bit patterns without spatial relationships.

\subsubsection{Quantum Circuit: BondFeatureMap}

The BondOrderMatrix was mapped to a parameterized quantum circuit using PennyLane (v0.32)\cite{pennylane}:

\textbf{(i) Initial encoding layer:}
\begin{equation}
|\psi_0\rangle = \bigotimes_{i=1}^{n} R_Y(\theta_i) |0\rangle, \quad \theta_i = \mathbf{M}_{ii}
\end{equation}

\textbf{(ii) Entangling layer:}
\begin{equation}
U_{\text{ent}} = \prod_{\langle i,j \rangle} R_{ZZ}(\phi_{ij}), \quad \phi_{ij} = \mathbf{M}_{ij}
\end{equation}

where $R_{ZZ}(\phi) = \exp(-i \phi Z \otimes Z / 2)$ creates two-qubit entanglement encoding off-diagonal matrix elements.

The circuit depth was set to 1 layer (encoding only) for the proof-of-concept Phase 1 study. The number of qubits was determined by the number of heavy atoms (C, N, O, S, P) in each molecule, resulting in circuits ranging from 12 to 28 qubits for the 17-molecule panel.

\subsubsection{Quantum Kernel Computation}

The quantum kernel between molecules $A$ and $B$ was computed as:

\begin{equation}
K(A, B) = |\langle 0 | U_B^\dagger U_A | 0 \rangle|^2
\end{equation}

where $U_A$ and $U_B$ are the BondFeatureMap circuits for molecules $A$ and $B$, respectively. The kernel represents the squared overlap of quantum states, measuring structural similarity in Hilbert space.

For the $n=\num{17}$ molecule dataset, the $n \times n$ kernel matrix $\mathbf{K}$ was computed using PennyLane's statevector simulator (\texttt{default.qubit} backend). Exploiting kernel symmetry ($K(A,B) = K(B,A)$), only $n(n+1)/2$ unique entries were evaluated. For $n=\num{17}$ molecules, this required \num{153} kernel evaluations.

Computation time per kernel entry was logged to assess scalability. At approximately \qty{5}{\minute} per kernel entry (28-qubit circuits, CPU single-threaded), total computation time was $\sim$\qty{13}{\hour}. Parallelization across 8 cores reduced wall-clock time to $\sim$\qty{1.6}{\hour}. All quantum simulations were performed on \PENDING{[hardware: CPU model, RAM]}.

\subsection{Classification with Quantum Kernel}

The quantum kernel matrix $\mathbf{K}$ was used as a precomputed kernel in scikit-learn's SVM classifier:

\begin{verbatim}
SVC(kernel='precomputed', C=1.0)
\end{verbatim}

Classification was performed using LOO-CV, matching the ECFP4 baseline protocol. For each fold, the $(n-1) \times (n-1)$ training kernel and $(n-1) \times 1$ test kernel were extracted from the full $n \times n$ matrix.

\subsection{Statistical Analysis}

\subsubsection{Performance Metrics}

Primary: AUC-ROC (accounts for class imbalance).

Secondary: Accuracy, precision, recall, F1-score, confusion matrix.

For the 15:2 inactive:active imbalance, we report balanced accuracy:
\begin{equation}
\text{Balanced Acc} = \frac{\text{Sensitivity} + \text{Specificity}}{2}
\end{equation}

\subsubsection{Quantum vs. Classical Comparison}

Performance difference was assessed using a paired $t$-test on AUC scores across the 17 LOO folds (one per molecule). Effect size was quantified using Cohen's $d$:

\begin{equation}
d = \frac{\text{AUC}_{\text{quantum}} - \text{AUC}_{\text{ECFP4}}}{\sigma_{\text{pooled}}}
\end{equation}

Statistical significance was declared at $\alpha = 0.05$. Following the honest-negative precedent established in Project 3\cite{sao2026p3}, we classify outcomes as:

\begin{itemize}
    \item \textbf{Advantage:} $\Delta\text{AUC} \geq 0.01$ and $p < 0.05$
    \item \textbf{Equivalence:} $|\Delta\text{AUC}| < 0.01$ or $p \geq 0.05$
    \item \textbf{Underperformance:} $\Delta\text{AUC} < -0.01$ and $p < 0.05$
\end{itemize}

\subsubsection{Kernel Target Alignment}

To assess whether the quantum kernel captures activity-relevant structure, we computed kernel target alignment (KTA)\cite{cristianini2002}:

\begin{equation}
\text{KTA} = \frac{\langle \mathbf{K}, \mathbf{yy}^\top \rangle_F}{\|\mathbf{K}\|_F \|\mathbf{yy}^\top\|_F}
\end{equation}

where $\mathbf{y}$ is the label vector and $\langle \cdot, \cdot \rangle_F$ is the Frobenius inner product. KTA ranges from $-1$ (perfect anti-correlation) to $+1$ (perfect alignment); values near 0 indicate no correlation between kernel similarity and label similarity.

\subsection{Computational Provenance}

All analyses were conducted using Python 3.10 in a reproducible conda environment (\texttt{malaria\_qml\_hybrid}). Key package versions: PennyLane 0.32.0, PyTorch 2.0.1, RDKit 2023.09.1, scikit-learn 1.3.0, NumPy 1.24.3, Pandas 2.0.3. Random seeds were fixed (seed=42) for all train/test splits. Quantum circuit parameters, kernel matrices, and SVM models were serialized and archived alongside results. Detailed provenance metadata, including computation timestamps and hardware specifications, are provided in the Supporting Information.

Code and data are available at \PENDING{[GitHub repository or Zenodo DOI upon acceptance]}.
